\subsubsection*{P1: static vs.\ dynamic assignment}
\inputif{q1_p1_static.tex}
\inputif{q1_p1_cyclic.tex}
\inputif{q1_p1_chunk.tex}
\inputif{q1_p1_dyn.tex}
\inputif{q1_p1_sched.tex}

\begin{itemize}
\item \textbf{Static and cyclic assignment depend strongly on the thread count}
  (Tables~\ref{tab:p1static}, \ref{tab:p1cyclic}): from $2^{16}$ to $2^{22}$
  threads the time drops from 9--14~s to about 1.5~s (3 repetitions per configuration). With few threads, each thread
  owns a fixed list of many operations that it must finish alone. The cost of an
  operation varies a lot (a few probes early on, thousands of probes
  when the table is nearly full), so the threads that happen to own expensive
  operations decide the run time while the others have finished and sit idle (load
  imbalance). Static assignment is only good when every thread gets a single
  operation ($2^{22}$ threads $=$ number of operations). Then the hardware
  block scheduler does the balancing: as soon as a block finishes, a new one
  starts in its place.
\item \textbf{Dynamic assignment with small chunks is insensitive to the thread
  count} (Table~\ref{tab:p1dyn}: 1.37--1.53~s for every configuration from
  $2^{16}$ to $2^{22}$ threads), because a thread that finishes early simply takes
  more work. Even $2^{16}$ threads (only 2.7$\times$ the
  \devmaxresidentthreads{} resident threads) are enough.
\item \textbf{The optimal chunk size is 1} (Table~\ref{tab:p1chunk}), and the time
  rises steadily with the chunk size (about 1.4~s at chunk 1, about 17~s at chunk 256).
  We expected larger chunks to help, as they do on the CPU, where they amortise the
  cost of the shared counter, but here the opposite happened. We think there are two
  reasons. (i) The \texttt{atomicAdd} on one counter is
  cheap compared with a single operation here (hundreds to thousands of
  dependent memory accesses), so there is nothing to amortise. (ii) Operations become more expensive the
  fuller the table is, so the last chunks handed out are the most
  expensive ones. A large last chunk makes one thread execute many of the most expensive
  operations \emph{sequentially} at the end, while the rest of the GPU is idle. Chunk 1
  minimises this tail.
\item \textbf{Block size} has only a second-order effect for P1 (Table~\ref{tab:p1dyn}):
  a thread does not cooperate with the other threads of its block, so the block
  size only affects how many blocks fit on an SM. 1024 was marginally best and
  is used below.
\item Best P1 configurations (Table~\ref{tab:p1sched}): dynamic (chunk \ponechunk)
  \ponedynamicms~ms, static \ponestaticms~ms, cyclic \ponecyclicms~ms. With one operation per
  thread the three schemes are nearly identical; dynamic wins slightly and, more
  importantly, does not need the thread count to match the operation count.
\end{itemize}

\subsubsection*{P2: thread-group size}
\inputif{q1_p2_group.tex}
\inputif{q1_p2_tune.tex}

\begin{itemize}
\item \textbf{The larger the group, the faster} (Table~\ref{tab:p2group}); $g=32$ is
  optimal at every thread count. A group of $g$ lanes examines $g$ consecutive
  slots with one coalesced load, so a probe sequence of length $L$ needs about $L/g$
  load instructions, ballots and loop iterations instead of $L$. With $g=32$ each
  probe step consumes exactly one 128-byte line, i.e.\ every byte that is
  fetched is used. Going from $g=16$ to $g=32$ still saves about a third of the time. Probe
  sequences in a full table are long (hundreds to thousands of slots), so
  even 32-slot windows are mostly useful.
\item \textbf{$g=1$ is about 2.5$\times$ slower than P1}, although it also uses one
  thread per operation. Each lane of a warp is its own group and runs its own probe loop
  containing \texttt{\_\_ballot\_sync}/\texttt{\_\_shfl\_sync} with a one-lane mask. Our
  explanation is that these synchronising intrinsics stop the compiler from running the
  32 independent loops of a warp as one predicated SIMD loop (as it does for P1), so
  the scheduler executes the lanes' loops as separate sub-warps (independent
  thread scheduling, lecture 6) and SIMD efficiency drops. We did not confirm this by looking at the generated code.
\item \textbf{Many more threads than can be resident help P2}: the best thread
  count is $2^{26}$, i.e.\ $2^{21}$ groups of 32 for $2^{22}$ operations, so each group
  gets only two operations. Like the static P1, P2 assigns operations to
  groups statically (cyclically). Our understanding is that making the work per group as
  small as possible lets the block scheduler balance the load, which removes most of
  the slow tail caused by the expensive late inserts. With $2^{20}$ threads (one group per 128 operations) P2
  is 4$\times$ slower than with $2^{26}$.
\item \textbf{Block size}: 128 is best (Table~\ref{tab:p2tune}). 1024-thread blocks
  are worse. Our explanation is that a block can only be replaced when \emph{all} of its 32 warps are
  done, so one warp stuck on a long probe sequence holds 1024 thread slots of the
  SM idle. Smaller blocks release their resources sooner. Below 128 the per-SM
  limit of \devmaxblockspersm{} blocks keeps the SM from being filled (16 blocks
  $\times$ 64 threads = 1024 $<$ \devmaxthreadspersm).
\end{itemize}

\subsubsection*{P2 with dynamic group assignment}
\inputif{q1_p2dyn_chunk.tex}
\inputif{q1_p2dyn_group.tex}

\begin{itemize}
\item \textbf{Dynamic assignment removes static P2's dependence on the thread
  count} (Table~\ref{tab:p2dynchunk}). With chunk 1 the time is 260--266~ms for
  every thread count from $2^{16}$ to $2^{24}$, whereas static P2 needs
  $2^{26}$ threads to reach 351~ms; with $2^{16}$ threads it takes 4.0~s, about
  15$\times$ longer than dynamic P2 with the same number of threads. Static P2 balanced the expensive late inserts only indirectly, by
  making the work per group so small that the block scheduler could balance it.
  A shared counter balances it directly: a group that finishes early takes the
  next operation.
\item \textbf{The best dynamic P2 (\ptwodynbestms~ms) takes 26\% less time than the best
  static P2 (\ptwobestms~ms)}, even though we tried static P2 with every thread
  count. We think the reason is that static P2 needs 524\,288 blocks to balance its
  load, and most of them start, find no work and exit, while dynamic P2 does the same
  balancing with a few thousand long-lived groups.
\item \textbf{Chunk 1 is optimal}, and the time grows steadily with the chunk
  size (about 260~ms at chunk 1, about 550~ms at chunk 16), for the same reason as
  in P1. One \texttt{atomicAdd} per group is negligible compared with an operation
  (tens to thousands of probes), while larger chunks put many of the most
  expensive operations, the last inserts into the full table, into one group's
  sequential tail.
\item \textbf{Group size 32 remains best} (Table~\ref{tab:p2dyngroup}). This is
  what we expected, since the coalescing advantage of wide windows should not depend
  on how operations are assigned to groups.
\end{itemize}

\subsubsection*{P1 vs.\ P2}
\inputif{q1_final.tex}

Table~\ref{tab:q1final} compares the best P1 (dynamic, chunk~\ponechunk, block size
\ponebs) with the best static P2 (group~\ptwogroup, block size \ptwobs) and the
best dynamic P2 (chunk~\ptwodynchunk, group~\ptwodyngroup), each at its best
thread count, for the five operation mixes and $2^{20}$, $2^{22}$ and $2^{24}$
operations. \qoneallverified{} runs passed verification. Dynamic P2 is the fastest
program in every case: \qonedynspeedupmin--\qonedynspeedupmax$\times$ faster than P1
and \qonedynvsstaticmin--\qonedynvsstaticmax$\times$ faster than static P2.
\begin{itemize}
\item \textbf{P2 is faster in every case, by \qonespeedupmin--\qonespeedupmax$\times$.}
  We think the difference comes mainly from memory-access coalescing and cooperation. P2 finds the target slot of an
  operation in $L/32$ steps, each a fully used 128-byte access, while a P1 thread walks the same
  $L$ slots one dependent load at a time. In addition, the 32 lanes of a P1 warp probe 32
  unrelated parts of the table, so every warp load touches up to 32 different lines
  (memory divergence: the warp waits for the slowest lane). P1's loads still
  hit in L1 for consecutive slots of the same thread, which is probably why the gap
  is ``only'' 4--6$\times$ and not 32$\times$.
\item \textbf{Time grows much faster than the operation count}: for 100\% inserts,
  $2^{20}\to2^{22}\to2^{24}$ operations take 0.15~s $\to$ 1.4~s $\to$ 67~s (P1)
  and 31~ms $\to$ 0.35~s $\to$ 14.9~s (static P2), although the operation count only
  grows by 4$\times$ each time. The table is sized to the number of distinct keys, so
  it ends 100\% full. For linear probing the total cost of filling a table
  completely grows like $n^{3/2}$: the last few percent of inserts have to walk
  very large clusters that keep merging. At $2^{24}$ the growth is even steeper,
  and we think the effect described next is part of the reason.
\item \textbf{Modulo bias of $h(k)=k \bmod n$ seems to matter in a full table.} Keys are uniform in
  $[0,N)$ with $N=2^{30}$. If $N/n$ is not an integer, $k \bmod n$ is not uniform: with
  $N/n = q + f$ the first $f\cdot n$ slots receive keys with probability proportional to
  $q+1$, the rest proportional to $q$. In a half-empty table such a small imbalance
  ($1/q$) probably does not matter. In a 100\% full table, however, every surplus key of the
  over-loaded region must overflow into the under-loaded one, which we expected to
  create one very large cluster that all late inserts walk through. We tested this
  with P2 on $2^{23}$ distinct insert keys:
  \begin{center}\small
  \begin{tabular}{rrr}\toprule
  table size $n$ & $N/n$ & time (ms)\\\midrule
  8\,388\,608 ($=2^{23}$) & 128.000 & 667\\
  8\,388\,000 & 128.009 & 695\\
  8\,356\,092 & 128.5 & 2374\\\bottomrule
  \end{tabular}
  \end{center}
  The same keys run $3.5\times$ slower when $N/n$ has fractional part $\frac12$ (half the
  table over-loaded by $1/128$). For our workloads $N/n\approx1024.5$, $256.5$ and $64.5$
  for $2^{20}$, $2^{22}$ and $2^{24}$ operations: the relative imbalance $1/q$ grows
  from 0.1\% to 1.6\%, which we think explains at least part of the extra jump at
  $2^{24}$.
\item \textbf{Effect of the operation mix.} Most of the cost comes from the inserts,
  because they fill the table, but the deletes and searches are not free either:
  \begin{itemize}
  \item 100/0/0 (all inserts) is the most expensive mix at $2^{22}$ and $2^{24}$, but
    not at $2^{20}$, where the two 80\% mixes take longer for all three programs
    (28--30\% for P1, 9--18\% for static P2 and 65--66\% for dynamic P2). We first thought this was due to the particular key sets, so we
    generated eight random key sets for each of these mixes at $2^{20}$ (P2 dynamic,
    $2^{20}$ threads). With every one of the eight seeds both 80\% mixes were slower than
    100/0/0 (25.9--31.0~ms against 23.3--28.9~ms), so the effect is real. Timing the insert phase alone (by giving the
    programs only the insert count), the inserts of an 80\% list take 17.6--19.3~ms
    against 25.2--28.0~ms for a 100\% list, a ratio of about 0.7, close to the
    $0.8^{3/2}\approx0.72$ that the $n^{3/2}$ growth predicts. The extra 11--12~ms
    come from the 20\% deletes or searches: in a full table each of them walks from
    the key's home slot to the key, which is a long way on average. We measured the
    same split at the larger sizes (P2 dynamic, best thread count). The deletes or
    searches cost 58--69\%, 41--49\% and 26--27\% of the time of the insert phase
    at $2^{20}$, $2^{22}$ and $2^{24}$, so their share shrinks as the table grows. At
    the same time the inserts of an 80\% list take only 0.43 ($2^{22}$) and 0.51
    ($2^{24}$) of the time of a 100\% list, much less than the 0.72 of the $n^{3/2}$
    rule. All these lists have $N/n$ with a fractional part of about 0.5, so we think
    the reason is the modulo bias described above: the large cluster it creates grows
    quickly with $n$ and makes the larger 100\% table relatively more expensive. So
    at $2^{22}$ and $2^{24}$ the extra inserts of the 100\% list cost more than the
    deletes or searches of the 80\% lists, and 100/0/0 is the most expensive mix.
  \item 50/50/0 is the cheapest at $2^{24}$, which surprised us at first. Every delete
    needs its own distinct key, so in this mix all inserted keys are distinct and
    $n$ is \emph{exactly} $2^{23}$. Then $N/n=128$ is an integer, there is no modulo
    bias, and the insert phase costs only 0.67~s instead of 2.39~s for the equally
    long insert phase of 50/0/50, whose $n$ is slightly smaller than $2^{23}$
    because of duplicate keys (we timed the insert phases separately). Each
    delete is a positive search plus one CAS.
  \item 50/0/50 and 50/50/0 are otherwise similar (at $2^{20}$ and $2^{22}$, where
    neither $n$ is a power of two): both do the same inserts, and a positive search
    and a delete both walk to the key and stop.
  \item 80/20/0 and 80/0/20 cost almost the same, for the same reason; a delete
    adds only one CAS.
  \end{itemize}
\item \textbf{The speedup of P2 is somewhat smaller at $2^{24}$ (3.9--4.5$\times$) than at
  $2^{20}$ (4.8--5.8$\times$).} We are not completely sure why. One possible
  explanation is that at the largest size both programs spend most of their time
  walking very long clusters near the end, where P1's sequential walk also uses whole
  cache lines through L1, and both programs compete for the few remaining empty
  slots.
\end{itemize}

\subsubsection*{Robustness across random key sets}
\inputif{q1_seeds.tex}

To check that the conclusions do not depend on one particular random operation
list, we regenerated the $2^{22}$-operation lists of all five mixes with two
further seeds and ran the best configuration of each program on them
(Table~\ref{tab:q1seeds}).
\begin{itemize}
\item \textbf{Absolute times vary by up to about 50\% between key sets} of the same
  mix (6--47\% depending on mix and program, e.g.\ 50/50/0: 871--1181~ms for P1). Our first guess was the modulo bias
  above, but the data does not support it. For every mix the three seeds have almost the
  same $N/n$ (fractional part $\approx0.49$), and in the 50/50/0 mix $n=2^{21}$
  exactly, so there is no bias at all, yet that mix varies by 35\%. We think the
  variation comes from the 100\% full linearly probed table itself: its cost is
  dominated by the few largest clusters, which the last inserts must walk, and the
  size of those clusters depends strongly on the exact keys (the cost seems to be
  heavy-tailed).
\item \textbf{The comparisons are stable}: P1/P2 lies in
  \qoneseedratioonemin--\qoneseedratioonemax$\times$ and static/dynamic P2 in
  \qoneseedratiotwomin--\qoneseedratiotwomax$\times$ across all seeds, and the
  ranking dynamic P2 $<$ static P2 $<$ P1 holds in all \qoneseedsameorder{} of
  the 15 seed/mix combinations. So we think the conclusions of this section do not
  depend on the particular random lists. Absolute times of a single list should
  be read with this key-set variation (up to about 50\%) in mind; the repetitions
  (3--5 per point) only remove timing noise, not this variation.
\end{itemize}

